\section{Local recognition of the full period group}
\label{sec:repeat}
A repeated shape need not identify a period. Two equal disks in one
periodic cell can be exchanged by a translation which fails to preserve
all the other disks. We therefore test a candidate translation on a
whole central patch, rather than identify bodies by their individual
shapes. The bounded period presentation makes this a finite test of a
global property. This removes the translation-type separation hypothesis
from exact fixed-aperture rigidity, including for nonprimitive supplied
presentations. A quantitative patch condition, which is automatic for
each fixed table but need not hold uniformly near a symmetry increase,
gives the finite-confidence version.

\subsection{The class and its observation}
\begin{definition}[Periodic class with repeated shapes]\label{def:repeatclass}
Let $\mathcal B$ satisfy the geometric and smoothness bounds of
Definition~\ref{def:aperture-prior}, except for the type condition
\eqref{eq:aperture-type}. More explicitly, its members are nonempty unions
\[
 \mathcal O=\bigcup_{i=1}^{r}(C_i+\Lambda_0),\qquad
 \Lambda_0=L\Z^2,\quad 1\le r\le r_0,
\]
with exactly $r$ physical component orbits modulo $\Lambda_0$,
$\|L\|,\|L^{-1}\|\le L_0$, and $\covol\Lambda_0\ge V_0>0$.
Different physical bodies have distance at least $d_0>0$; their diameters
are at most $D_0$, curvatures belong to $[\kappa_-,\kappa_+]$, and
Steiner-centered support functions have $C^{6,\beta}$ norm at most $M_6$.
The presentation is not supplied to the observer and need not be
primitive. Identical shapes, arbitrary body symmetries and repeated
copies within its cell are allowed. No shape-separation, asymmetry,
clear-bridge, area-normalization or return-window assumption is made.
\end{definition}

Write $\mathcal C$ for the physical components, $z_C=s(C)$ for the
Steiner point, and $p_C$ for the centered support function in the common
laboratory directions. The unknown intrinsic period group is
\[
       \Pi(\mathcal O)=\{v\in\R^2:\mathcal O+v=\mathcal O\}.
\]
A multiple of a period cell is not an additional unknown: the conclusion
concerns $\Pi(\mathcal O)$ itself and its component orbits. Put
\begin{equation}\label{eq:repeatwindow}
 B=1+L_0,\qquad \tau=\min(1,d_0/8),\qquad
 R_*=12B+2D_0+3,\qquad Q_a=[-a,a]^2.
\end{equation}
Each attempt selects a uniform position in $Q_{R_*}$ and an independent
uniform direction. A solid start, no impact before $\tau$, or a first
impact outside $Q_{11B+D_0+1}$ is a failure. Otherwise only the first
impact position is retained. Denote this position-and-failure law by
$\mathsf F^*_{\mathcal O}$. Its normalization uses the known launch
square. No collision angle, later collision, elapsed time within the
cutoff, boundary coordinate, type label or exact cell area is observed.
Finite observations have certified position error $\varepsilon_x$.
These remain controlled independent launches in a common laboratory
frame; they are not observations of one passive orbit.

\begin{theorem}[Rigidity with arbitrary repeated shapes]
\label{thm:repeat-exact}
On $\mathcal B$, the support of the single first-impact position law
$\mathsf F^*_{\mathcal O}$ determines the whole periodic obstacle union,
its full translation lattice $\Pi(\mathcal O)$, the number $r_*$ of
component orbits under that lattice, and the free area of its primitive
cell. No distinction between the individual obstacle shapes is required.
The determination is in the laboratory frame; forgetting that frame
leaves the usual common Euclidean isometry and lattice-basis ambiguities.
In particular, a bound on some period presentation suffices: neither a
primitive presentation nor a primitive-cell multiplicity is supplied.
\end{theorem}

The exact assertion is a finite-patch determination theorem within a
class already assumed periodic. It does not deduce crystallinity from
local congruence alone, and it does not identify the position law with a
marked length spectrum. Its proof uses no analytic continuation or
infinite contact-jet inverse.

\subsection{A finite certificate for a translation period}
Set
\begin{equation}\label{eq:repeatmetric}
 d_*(C,D)=\|z_C-z_D\|_\infty+\|p_C-p_D\|_\infty.
\end{equation}
This is a metric on translated convex bodies in laboratory coordinates.
It does not minimize over rotations. The Steiner-point formula implies
$\|z_C-z_D\|_\infty\le2d_H(C,D)$ and hence
$d_*(C,D)\le5d_H(C,D)$; conversely $d_H(C,D)\le\sqrt2d_*(C,D)$.
For any bounded $v$ define the central defect
\begin{equation}\label{eq:repeatdefect}
 \Delta_{\mathcal O}(v)=
 \max_{C:\ z_C\in Q_B}\ \max_{\epsilon\in\{-1,1\}}
          \inf_{D\in\mathcal C}d_*(C+\epsilon v,D).
\end{equation}
The central set is finite and nonempty. The infima are attained, since
centers far from the bounded set of translated sources cannot minimize.

\begin{lemma}[Local-to-global period test]\label{lem:period-local}
For every $v$, $\Delta_{\mathcal O}(v)=0$ if and only if
$v\in\Pi(\mathcal O)$. If $\|v\|_\infty\le8B$, the zero test uses
only components with centers in $Q_B$ and $Q_{9B}$. All the component
orbits modulo the unspecified $\Lambda_0$ are represented among the
central sources.
\end{lemma}
\begin{proof}
For any component center choose a lattice translate in
$L[-1/2,1/2]^2$. Its Euclidean norm is at most $L_0/\sqrt2<B$, so
there is a central representative of every $\Lambda_0$ orbit. If the
defect vanishes, each such representative $C$ satisfies
$C+v\in\mathcal C$ and $C-v\in\mathcal C$. Translating these statements
by $\Lambda_0$ gives $\mathcal O+v\subset\mathcal O$ and
$\mathcal O-v\subset\mathcal O$, hence equality. The converse follows
from the definition of a period. A component equal to $C\pm v$ has
center $z_C\pm v$, in $Q_{9B}$ under the stated bound. A zero test thus
requires only that finite set; it does not have to compute a distant
minimizer of a positive defect.
\end{proof}

\begin{lemma}[Full periods from a protected finite list]
\label{lem:repeatgroup}
The group $\Pi=\Pi(\mathcal O)$ is a lattice containing $\Lambda_0$,
with
\begin{equation}\label{eq:repeatindex}
              [\Pi:\Lambda_0]\le r\le r_0,
       \qquad \covol\Pi\ge V_*:=V_0/r_0.
\end{equation}
Fix any root component with center in $Q_{2B}$. Among the differences
$z_D-z_{C_*}$ for targets with centers in $Q_{4B}$, retain those which
pass the local period test. Their integer span is exactly $\Pi$.
\end{lemma}
\begin{proof}
Fix a component $C_*$. Map a period modulo $\Lambda_0$ to the
$\Lambda_0$-component orbit of $C_*+v$. This map is injective: two periods
with the same image differ by a vector of $\Lambda_0$, since a compact
body has no nonzero translational symmetry. Thus the quotient group
has at most $r$ elements. A finite extension of a rank-two lattice is
a rank-two lattice, and covolumes give \eqref{eq:repeatindex}.
The action of this quotient on component orbits is free, so its order
also divides $r$ and $r_*=r/[\Pi:\Lambda_0]$.

Every retained difference is a period by Lemma~\ref{lem:period-local}.
Conversely, the two columns of $L$ have norm at most $L_0<B$.
Each coset of $\Pi/\Lambda_0$ has a representative in
$L[-1/2,1/2]^2$, also of norm less than $B$. The root translated by
any of these vectors has center in $Q_{3B}$, so it is among the
candidate targets. The accepted list therefore contains generators of
$\Lambda_0$ and representatives of every coset of $\Pi/\Lambda_0$.
They generate $\Pi$. The gap from $Q_{3B}$ to the target cutoff $Q_{4B}$
is deliberately retained for the noisy construction.
\end{proof}

\begin{proof}[Proof of Theorem~\ref{thm:repeat-exact}]
The first-hit flux calculation in \eqref{eq:aperture-flux}, with $R_*$
in place of $R$, applies to every component meeting $Q_{11B}$.
Its whole boundary and a short exterior launch collar lie inside the
observation and launch squares. In particular every open boundary arc
has positive observed mass. The support of the observed measure has
no points off the physical boundaries. Physical separation therefore
identifies all complete components in the inner patch from this support.

There is no need to infer completeness by looking for a closed-looking
fragment. Take the convex hull of each connected support component and
keep those whose Steiner points lie in $Q_{10B}$. A partial outer
component has its hull inside its true body. If its hull center meets
that cutoff, its body meets $Q_{10B}$, so its whole boundary is observed
and belongs to the same connected component after all. Thus every kept
hull is a complete physical body. The same reasoning will apply to
finite clouds on their simultaneous coverage event.

Choose a root in $Q_{2B}$ and the targets in $Q_{4B}$. Test both
translations of every central body. Candidate vectors have sup norm at
most $6B$, and all matching targets lie inside the available $Q_{10B}$
patch. Lemmas~\ref{lem:period-local} and \ref{lem:repeatgroup} identify
the full period lattice. One may generate a basis from an independent
pair by the finite rational and Hermite procedure in
\eqref{eq:aperture-Q}--\eqref{eq:aperture-lattice}, now using $V_*$.
No area equation is used to decide which displacements are periods.

Take all bodies with centers in $Q_{2B}$ and quotient them by $\Pi$.
The central patch contains every $\Lambda_0$ orbit and hence every
$\Pi$ orbit. Two bodies are in the same $\Pi$ orbit precisely when
their center difference is a period: a period sends the first body to
some body with the second center, and two disjoint convex bodies cannot
have a common Steiner point. The same finite test decides this relation.
One representative from each class and the recovered lattice determine
$\mathcal O$ everywhere. Finally
\begin{equation}\label{eq:repeat-area}
           A_*=\covol\Pi-\sum_{j=1}^{r_*}\area(C_j)
\end{equation}
is the free area of a primitive cell. A multiple-cell presentation of
the same union produces the same answer.
\end{proof}

\subsection{Finite confidence without individual shape labels}
For exact recovery no positive mismatch bound was assumed. Uniformly
correct decisions from noisy data require a condition on nonperiods,
not on whether two individual shapes resemble each other.

\begin{definition}[Patch mismatch margin]\label{def:repeatmargin}
Let
\[
 \mathcal V_{\mathcal O}
   =\{z_D-z_C:\ z_C\in Q_{3B},\ z_D\in Q_{5B},\ C,D\in\mathcal C\}.
\]
The class $\mathcal B_\eta$, $\eta>0$, consists of those tables in
$\mathcal B$ for which
\begin{equation}\label{eq:repeatmargin}
 v\in\mathcal V_{\mathcal O}\setminus\Pi(\mathcal O)
                \quad\Longrightarrow\quad\Delta_{\mathcal O}(v)\ge\eta.
\end{equation}
There is no restriction on the number of congruent bodies or on their
symmetry groups. For each fixed table the candidate set is finite and
Lemma~\ref{lem:period-local} implies some positive margin. The bound is
therefore pointwise automatic, though its uniform numerical value is
additional information in a finite-confidence theorem.
\end{definition}

\begin{theorem}[Finite first-impact recovery of repeated motifs]
\label{thm:repeat-finite}
Fix the numerical bounds of $\mathcal B_\eta$. There are
$C,c,\nu_0>0$ depending only on those bounds such that, for
$0<\nu<\nu_0$ and $0<\delta<1/4$, at most
\begin{equation}\label{eq:repeatcost}
       C\nu^{-3/4}\log\frac{C}{\nu\delta}
\end{equation}
independent launches from the fixed square \eqref{eq:repeatwindow},
with position error $\varepsilon_x\le c\nu^{3/2}$, suffice with
probability at least $1-\delta$ for the following conclusions.
The procedure recovers $r_*$ and the exact rational relations of an
accepted generating list for the full period lattice. It gives a finite
presentation of the entire periodic table with $C^2$ support error and
period-basis error at most $C\nu$, and $|\widehat A_*-A_*|\le C\nu$.
The metric is the finite-presentation metric following
Theorem~\ref{thm:aperture-main}, with component orbits under $\Pi$,
not equivalence classes of shapes, as its indices.

The construction uses a finite cloud, finite patch tests and bounded
rational searches. It does not receive a shape dictionary, a primitive
cell, an incidence graph, integer period labels, return histograms or
adaptive recentered launches. All failures are counted in
\eqref{eq:repeatcost}. The patch margin, lattice-presentation bound,
finite smoothness, curvature, separation and localization accuracy are
explicit hypotheses; the launch bound does not include arithmetic,
localization bit cost or apparatus travel and setup.
\end{theorem}

\begin{proof}
\emph{Simultaneous geometric coverage.}
A separation packing bounds the number of bodies meeting $Q_{11B}$
by a prior-dependent constant $M_*$. For an arc $I$ on any such body,
\[
 \mathsf F^*_{\mathcal O}(I)\ge
                 \frac{\tau|I|}{8\pi(4R_*^2)}.
\]
Indeed the launch tube $q=x-rv$, $\tau/4\le r\le\tau/2$, with directions
within $\pi/3$ of the inward normal, has first-hit Jacobian $\cos\phi$.
Convexity excludes an earlier hit on the same body, and $r<d_0$
excludes another body. Thus this is an unconditional probability per
attempt, not uniform sampling of an unknown boundary. Divide each
boundary into arcs of lengths between $q/2$ and $q$. A union bound gives
complete $q$-coverage with probability at least $1-\delta$ when
\begin{equation}\label{eq:repeatcoupon}
 N\ge\frac{16\pi(4R_*^2)}{\tau q}
             \log\frac{2M_*(1+\pi D_0/q)}{\delta}.
\end{equation}
On this event, join measured points at distances less than $d_0/3$,
and take their cluster hulls $P_C$. With sufficiently small $q$ and
$\varepsilon_x$, no distinct bodies merge and every relevant complete
boundary gives one connected cluster. As in \eqref{eq:aperture-hull},
\begin{equation}\label{eq:repeathull}
         d_H(P_C,C)\le e:=\kappa_+q^2/2+\varepsilon_x.
\end{equation}
Keep hulls with raw Steiner center in $Q_{10B}$. Even an initially
partial cluster has its hull in the $\varepsilon_x$-neighborhood of one
true body. A kept hull therefore belongs to a body meeting $Q_{11B}$;
coverage and connectivity make that cluster complete. All subsequent
tests use only these certified-by-coverage complete hulls.

\emph{Uniform control of each finite patch test.}
For the moment compute centers and support suprema exactly from the
finite polygons. These operations admit certified finite numerical
approximation, accounted for below. The raw center error is at most
$2e$ and the centered support error at most $3e$. Select sources by
$\|s(P_C)\|_\infty\le2B$, a root from them, and targets for the root by
$\|s(P_D)\|_\infty\le4B$. For each root-target difference
$\widehat v=s(P_D)-s(P_{C_*})$ calculate
\begin{equation}\label{eq:repeatmeasured}
 \widehat\Delta(\widehat v)=
 \max_{C:\ s(P_C)\in Q_{2B}}\ \max_{\epsilon\in\{-1,1\}}
   \min_{E:\ s(P_E)\in Q_{10B}}d_*(P_C+\epsilon\widehat v,P_E).
\end{equation}
The paired true vector $v=z_D-z_{C_*}$ has error at most $4e$.
For each fixed source-target comparison, the center term changes by
at most $8e$ and the centered-support term by at most $6e$.
Their sum is at most $14e$. If $v$ is a true period, every measured
source has its true translated partners inside $Q_{8B+8e}$; their raw
centers belong to $Q_{10B}$ for small $e$. Hence
$\widehat\Delta(\widehat v)\le14e$.

If $v$ is not a period, the true root and target centers lie in
$Q_{3B}$ and $Q_{5B}$, so \eqref{eq:repeatmargin} applies. Every witness
source with true center in $Q_B$ is selected by the $Q_{2B}$ cutoff.
Restricting possible targets cannot reduce the true infimum. It follows
that $\widehat\Delta(\widehat v)\ge\eta-14e$. Allow a further $6e$ for
certified center and support computations and use the threshold $\eta/2$.
For $20e<\eta/4$, accepted vectors are exactly the true periods among
the candidates. In particular, neither an optional outer cluster nor
a false individual shape match can introduce a spurious period.

The same test is used on every difference between two source centers.
Their true endpoints are in $Q_{3B}\subset Q_{5B}$, and all translated
source targets remain strictly inside $Q_{10B}$. The preceding bounds
still hold. Acceptance is consequently the exact equivalence relation
of belonging to the same $\Pi$-component orbit, and determines $r_*$.
No comparison of arbitrarily close independent shape classes is made.

\emph{Protected generators and exact arithmetic.}
A central representative of each $\Lambda_0$ orbit has true center
strictly inside $Q_B$, so a measured root in $Q_{2B}$ exists. Its
translates by the two columns of $L$ and by the bounded coset
representatives in Lemma~\ref{lem:repeatgroup} all have measured centers
in $Q_{3B+4e}$ and are retained in $Q_{4B}$. Thus the accepted list
spans exactly $\Pi$, even if other candidates cross selection cutoffs.
Every accepted vector has Euclidean norm at most $U=12B$ for small $e$.
Every nonzero determinant of two periods is an integer multiple of
$\covol\Pi\ge V_*$, whereas its perturbation is $O(Ue)$.
Independence can therefore be detected at the threshold $V_*/2$.
Choose an estimated maximal-determinant independent pair and call its
true matrix $D$. Set
\[
                         Q=\max(1,\lceil U^2/V_*\rceil).
\]
The lattice $D\Z^2$ has index $n\le Q$ in $\Pi$. Each coordinate of
$D^{-1}v_j$ has reduced denominator dividing $n$, and bounded numerator.
The coordinate errors are at most $C_*e$ by the inverse-matrix formula.
Distinct rationals of denominator at most $Q$ are separated by
$Q^{-2}$. Once $C_*e<1/(3Q^2)$, the nearest admissible rational is
unique and is the exact coordinate. Thus the rational decisions are
made only after geometric candidates have been certified as periods.

Let $q_*$ be the common denominator and let $H$ be a column Hermite
basis of the integer span of $q_*I_2$ and $q_*D^{-1}v_j$. Then
\begin{equation}\label{eq:repeatbasis}
              \Pi=(DH/q_*)\Z^2,\qquad
              \widehat L_*=\widehat D H/q_*.
\end{equation}
Here $q_*\le Q$, and the Hermite entries are bounded by $Q^2$ because
the integer group contains $q_*\Z^2$. The basis error is $O(e)$.
As usual, matched bases rather than continuously chosen reduced bases
are used. Pick one raw hull per recovered component orbit. Its area
error is $O(e)$ by the parallel-body bound, and \eqref{eq:repeat-area}
then gives $|\widehat A_*-A_*|=O(e)$.

\emph{Smooth geometry and the resource bound.}
Use the same compensated kernel as \eqref{eq:aperture-smoothing}.
On centered supports its error obeys
\[
 \|K_h*p_{P_C}-p_C\|_{C^2}
                 \le2C_\varphi e h^{-2}+B_6h^4.
\]
The common $C^{6,\beta}$ prior supplies $B_6$. Choose
$h\asymp\nu^{1/4}$ and $e\asymp\nu^{3/2}$ with small fixed constants,
and take $q\asymp\nu^{3/4}$ and
$\varepsilon_x\le c\nu^{3/2}$ in \eqref{eq:repeathull}.
Then the support error is $O(\nu)$, and
$\widetilde p+\widetilde p''>1/(2\kappa_+)$ for small $\nu$;
the signed smoothing therefore defines a strictly convex body.
Centers are restored after smoothing. Numerical integrals of the
piecewise polygonal support, its supremum differences and its Steiner
point have finite certified enclosures: an angular mesh and the uniform
support Lipschitz bound control the suprema, and the explicit kernel
controls quadrature. Reduce these errors below $e$ in the patch test and
below a fixed fraction of $\nu$ in the smoothed support.
Decrease $\nu_0$ to enforce every displayed geometric, mismatch and
rational threshold. Equation~\eqref{eq:repeatcoupon} gives
\eqref{eq:repeatcost}. All decisions hold on the one simultaneous
coverage event, including all choices made from that cloud. No
additional independence assumption for the selected root is used.
Finally, only bounded integer translates of the reconstructed motif
can enter a bounded fundamental region, since the matched basis and its
inverse are bounded. The original positive separation therefore persists
for sufficiently small perturbations. This proves the finite-presentation
and whole-table conclusions.
\end{proof}

\subsection{Pointwise recovery and the role of the margin}
\begin{corollary}[Unknown margin: eventual recovery, not premature certification]
\label{cor:repeatpointwise}
For every fixed table in $\mathcal B$, fresh dyadic batches with target
errors $e_j\downarrow0$, summable coverage failure probabilities, and
patch threshold $\sqrt{e_j}$ eventually recover the correct full-period
relations and component-orbit count almost surely. No uniform finite
index at which those discrete decisions become correct is asserted
without a numerical patch margin.
\end{corollary}
\begin{proof}
The relevant bounded physical pairs form a finite set. Every nonperiod
candidate has positive defect by Lemma~\ref{lem:period-local}, so the
minimum over all such pairs is positive for this fixed table. On a
coverage event, period defects are at most $20e_j$ and nonperiod defects
are at least that minimum minus $20e_j$. The threshold $\sqrt{e_j}$
separates them eventually. All protected generators are present, and the
rational thresholds eventually hold. The first Borel--Cantelli lemma
makes only finitely many coverage events fail. Recompute decisions from
each fresh cloud rather than permanently keeping an earlier class label.
The geometric estimates then converge with the stated eventual discrete
recovery. This is consistency, not an observable finite stopping certificate
on the union of all margin classes.
\end{proof}

\begin{proposition}[Repeated disks and a vanishing patch margin]
\label{prop:repeatjump}
The exact theorem includes repeated equal disks not covered by the
translation-type-separated theorem. Correct primitive multiplicity and
full lattice cannot be uniformly continuous under geometric perturbation
on all of $\mathcal B$ without a condition excluding increases of
translation symmetry.
\end{proposition}
\begin{proof}
Let $0\le t<1/8$ and place disks of radius $1/4$ at
\[
       (4k,4l),\qquad(4k+2+t,4l),\qquad k,l\in\Z.
\]
The diameter, curvature, smoothness, presentation and physical-separation
bounds are uniform. At $t=0$ the full period lattice is
$2\Z\times4\Z$ and there is one component orbit. For $t>0$ it is
$4\Z\times4\Z$ and there are two. To verify this, a period must send
center rows to center rows; horizontally it must preserve
$\{0,2+t\}$ modulo $4$. The nonzero exchange is possible precisely
when $2(2+t)$ is a multiple of $4$, which in this interval occurs only
at $t=0$.

The root-to-other-disk displacement $v=(2+t,0)$ matches one entire
body exactly, but for $t>0$ its patch defect is $2t$: the other body
would have to move to a center at $4+2t$, whereas the nearest such
center is $4$. The backward test gives the same discrepancy. Thus the
patch test rejects a false period despite complete equality of individual
shapes, and its margin tends to zero at the symmetry increase. The
physical unions converge on bounded windows at geometric error at most
$t$, while the primitive count changes and the covolume jumps from
$16$ to $8$. No modulus tending to zero for the full lattice and count
can hold across this family in the finite-presentation sense.
The statement concerns geometric or localization noise, not total
variation convergence of singular exact impact measures.
\end{proof}

\subsection{Relation to local symmetry criteria and the retained results}
Dolbilin's local theory~\cite{Dolbilin} gives criteria for crystalline
structure and global symmetry of Delone sets from local cluster
conditions. Our assumption is different: periodicity and a bound on
some lattice presentation are already given. Lemma~\ref{lem:period-local}
uses that bound to place representatives of every unknown component
orbit in one finite patch; it recovers the full group rather than
deducing periodicity of an arbitrary point set. Component centers alone
are insufficient when shapes differ, which is why the patch metric
includes complete centered supports. The proof is explicit and does not
invoke a general local-crystallinity theorem.

The launch-to-hull mechanism and compensated approximation remain
classical ingredients; compare \cite{Schneider,PSSW}. The additional
geometric result is the period test with repeated shapes, including its
primitive-presentation recovery and the discrete stability condition.
It is not a new general random-polytope rate. Boundary/lens rigidity
\cite{SUV} concerns unknown metrics and boundary data; obstacle
travelling-time rigidity~\cite{GNS23a,GNS23b} uses exterior entry, exit
and length data. Here the Euclidean metric and launch square are known,
and first impacts are localized inside a laboratory aperture. These data
are not asserted equivalent to a passive orbit, a marked spectrum or
collision counts.

All subsequent fixed-aperture, record-local, compensated-cloud, intrinsic
inverse and sequential results are retained with their earlier hypotheses.
The individual-shape gap in the next section remains a useful sufficient
condition for its simpler classifier; it is not retroactively imposed on
Theorem~\ref{thm:repeat-exact}. Likewise the mismatch example does not
weaken the earlier theorem: it identifies the proper distinction between
unconditional exact determination and uniform noisy recovery of a
symmetry-dependent discrete invariant. The earlier analytic, count-fiber,
moment and relative-law arguments remain supplied in the retained volumes.
